\section{Analytic realization and contraction}
\label{analytic:section}

We now realize the finite gluing data of Proposition~\ref{prop:jets} by
an actual complex surface.  The two prescribed vector fields will then
be lifted from their finite jets by a cohomology vanishing argument.
The realization uses only the finite list of genuine holomorphic flow
maps already constructed. The formal-functions theorem used afterwards
computes a cohomology group on that realized neighborhood.
Throughout this section, powers of an ideal denote ordinary ideal
powers, and cohomology is analytic sheaf cohomology. We use coherence
of analytic structure and ideal sheaves as in
\cite[Chapter~II, Theorems~3.19 and~4.29]{Demailly2012}.

\subsection{A Hausdorff neighborhood from the finite overlap maps}

\begin{lemma}\label{analytic:realization}
The overlap maps of Proposition~\ref{prop:jets} define, after shrinking
their domains, a smooth Hausdorff complex surface neighborhood $Y$ of
$C$.  Its conormal bundle is $L$, and hence
\[
 C^2=-10.
\]
Writing $\cI=\OO_Y(-C)$ and $\cF=\cT_Y(-\log C)$, the prescribed rows
define two global sections
\[
 \bar v_1,\bar v_2\in H^0(Y,\cF/\cI^9\cF).
\]
In divisor coordinates $(t,z)$, their local representatives have
logarithmic determinant $t^4(1+O(t))$ and ordinary determinant
$t^5(1+O(t))$.
\end{lemma}

\begin{proof}
Let $N=L^{-1}$ be the split normal bundle.  All changes to its
disk-to-outside identifications are finite compositions of flows of
holomorphic vector fields
\[
 D=t^k\bigl(b(z)t\partial_t+h(z)\partial_z\bigr),\qquad k\geq1,
\]
on punctured disk neighborhoods.  On a compact annulus in any such
overlap, these vector fields vanish along the zero section.
The local existence theorem for holomorphic differential equations,
applied uniformly on that compact annulus, therefore gives the
required time-one maps and their inverses after shrinking the fiber
radius.  There are only finitely many flows.  Shrinking successively
gives their compositions on neighborhoods of the same compact
annular bands.  Denote the resulting inside-to-outside maps by
$\Phi_j$, $1\leq j\leq6$.

Here is a quotient construction that also verifies the separation
property.  Choose mutually disjoint coordinate disks $D_j$ about the
six points, and in each choose radii
\[
 0<a_j<a'_j<b'_j<b_j
\]
so that $\Phi_j$ and its inverse are defined on neighborhoods of the
closed annulus $a_j\leq |z_j|\leq b_j$ in their respective split
bundles.  Their domains may be taken somewhat wider than these closed
annuli.  Retain a small open neighborhood $V_j$ of the zero section
over $|z_j|<b'_j$, and a small open neighborhood $V_0$ in $N$ over
\[
 C\setminus\bigcup_{j=1}^6\{|z_j|\leq a'_j\}.
\]
Choose the transverse thicknesses small enough that the maps are
defined throughout the retained parts of the closed annuli and
\begin{equation}\label{analytic:collar-bound}
 \bigl|z_j(\Phi_j(p))-z_j(p)\bigr|
 <\frac{a'_j-a_j}{2}
\end{equation}
there.  This is possible because each map restricts to the identity
on the zero section.  We also make the images of the retained annuli
lie over their respective disjoint disks $D_j$.

Set
\[
 A_j=\{p\in V_j:a_j<|z_j(p)|<b_j,\ \Phi_j(p)\in V_0\},
\]
and identify $p\in A_j$ with $\Phi_j(p)\in V_0$.  These are
biholomorphic identifications between open subsets.  A point of
$V_0$ is in at most one of their images, since those images lie over
disjoint disks.  Each equivalence class in
$V_0\sqcup V_1\sqcup\cdots\sqcup V_6$ thus has at most two elements.
In particular the equivalence relation consists exactly of the
diagonal, these six graphs, and their inverses; no further
identifications or triple-overlap equations occur.

Each graph is closed relative to $V_j\times V_0$.  Indeed, suppose
$p_\nu\in A_j$, $p_\nu\to p\in V_j$, and
$\Phi_j(p_\nu)\to q\in V_0$.  The source base radii have limit at
least $a_j$ and less than $b'_j<b_j$.  The flow is defined across the
closed annulus, so continuity gives $q=\Phi_j(p)$.  If the limiting
source radius were $a_j$, inequality~\eqref{analytic:collar-bound}
would give $|z_j(q)|<a'_j$, contrary to $q\in V_0$.  Thus
$a_j<|z_j(p)|<b_j$ and $p\in A_j$, as required.  Possible fiber
boundaries cause no difficulty because both limiting points are
assumed to belong to the retained open pieces.

The equivalence relation is therefore closed.  The quotient map is
open, since it glues open sets by biholomorphisms.  An open quotient
map with closed equivalence relation has Hausdorff target: two
inequivalent representatives have product neighborhoods disjoint
from the relation, and the images of these neighborhoods are
disjoint open sets.  Moreover, the quotient map restricts on every
$V_i$ to an injective open map.  The complex structures on these
pieces consequently give the quotient a smooth complex surface
structure.  The images of their zero sections give the original
compact curve $C$ with its original complex structure.

Every modifying flow fixes $C$ and induces the identity on its
conormal bundle.  In the split identification its coordinate
expansion has $t'=t+O(t^2)$, although its base coordinate may change
by $O(t)$.  Composing with the original split changes therefore
leaves $\cI/\cI^2\simeq L$.  This proves $C^2=-\deg L=-10$.

In the logarithmic basis $(t\partial_t,\partial_z)$ the prescribed
fields have coefficients $tA_i,tB_i$.  Rows with indices $0$ through
$7$ prescribe these coefficients through degree $8$ in $t$.
The matching in Proposition~\ref{prop:jets} is therefore precisely
matching modulo $\cI^9\cF$ on the actual overlap maps.  It gives the
claimed quotient sections.  Their logarithmic determinant is
\[
 t^2(A\wedge B)=t^4(1+O(t)).
\]
Replacing $t\partial_t$ by $\partial_t$ multiplies this determinant
by $t$, giving the stated ordinary determinant.
\end{proof}



\subsection{Vanishing and holomorphic lifting}

The finite jets now live on a Hausdorff surface. To lift them to
holomorphic fields, we prove that the first cohomology of
$\cI^9\cF$ vanishes after shrinking about $C$.
The logarithmic tangent sheaf $\cF$ is locally free of rank two.
Restriction to $C$ gives the exact sequence
\begin{equation}\label{analytic:log-sequence}
 0\longrightarrow\OO_C\longrightarrow\cF|_C
 \longrightarrow\cT_C\longrightarrow0.
\end{equation}
Here the first map sends $1$ to the class of the normal Euler field
$t\partial_t$, and the second is restriction of a tangent vector
field to $C$.  The Euler class is unchanged when the local equation
of $C$ is multiplied by a unit.  The sequence is thus intrinsic; a
splitting is unnecessary.

\begin{lemma}\label{analytic:graded-vanishing}
For every $n\geq9$,
\[
 H^1(C,\cF|_C\otimes L^n)=0.
\]
Consequently, for $\mathcal E=\cI^9\cF$ and every $r\geq1$,
\[
 H^1(Y,\mathcal E/\cI^r\mathcal E)=0.
\]
\end{lemma}

\begin{proof}
Since $K_C\simeq L^4$, tensoring
\eqref{analytic:log-sequence} by $L^n$ gives
\[
 0\longrightarrow L^n\longrightarrow\cF|_C\otimes L^n
 \longrightarrow L^{n-4}\longrightarrow0.
\]
Serre duality identifies the duals of the first cohomology groups of
the outer terms with
\[
 H^0(C,L^{4-n})\quad\hbox{and}\quad H^0(C,L^{8-n}),
\]
respectively.  Both vanish for $n\geq9$, as $\deg L>0$.
The first assertion follows from the cohomology sequence.

Because $C$ is a Cartier divisor and $\cF$ is locally free,
\[
 \cI^n\cF/\cI^{n+1}\cF\simeq\cF|_C\otimes L^n.
\]
The quotient
$\mathcal E/\cI^r\mathcal E=\cI^9\cF/\cI^{9+r}\cF$
has a finite filtration with precisely these graded terms for
$9\leq n<9+r$.  Their first cohomology vanishes by the first
assertion.  Successive cohomology sequences give the second
assertion; the quotient is supported on the compact curve $C$.
\end{proof}

\begin{lemma}\label{analytic:actual-lifting}
After shrinking $Y$ about $C$, the sections $\bar v_1,\bar v_2$ lift
to holomorphic sections $v_1,v_2$ of $\cF$.  They form a frame for
$\cT_{Y\setminus C}$.
\end{lemma}

\begin{proof}
The intersection matrix of $C$ is the negative definite matrix
$(-10)$.  Grauert's contraction criterion therefore contracts $C$
to a point $x$ of a normal complex surface:
\[
 \pi:Y\longrightarrow X,\qquad \pi^{-1}(x)=C,
 \qquad Y\setminus C\simeq X\setminus\{x\},
\]
where the fiber equality is set-theoretic, and
$\pi_*\OO_Y=\OO_X$.
We use the criterion in \cite[\S4.8(e), pp.~366--367]{Grauert1962}.
The contraction is proper and is a biholomorphism off the
exceptional curve.  It is the holomorphic reduction in that theorem,
whose target is normal when the source is normal; these properties
are part of the reduction and contraction
statements in \cite[\S2, p.~337 and Definition~3,
pp.~339--340]{Grauert1962}.  We choose a sufficiently small Stein
representative $X$ about $x$ and replace $Y$ by its inverse image.
In particular $X\setminus\{x\}$ is smooth.

We prove $R^1\pi_*\mathcal E=0$ for
$\mathcal E=\cI^9\cF$, keeping track of the two ideal filtrations.
Let $\mathfrak m_x$ be the ideal sheaf of $x$ in $X$ and put
$\mathcal J=\mathfrak m_x\OO_Y$.  Both $\mathcal J$ and $\cI$
vanish exactly on $C$, and $\cI$ is radical.  The analytic
Nullstellensatz~\cite[Chapter~II, Theorem~4.22]{Demailly2012}
and Noetherianity give an inclusion
$\cI^a\subseteq\mathcal J$ locally at each point of $C$ for some
positive integer $a$.  A finite cover of the compact curve gives a
single exponent.  After shrinking the representative,
\begin{equation}\label{analytic:cofinal}
 \cI^a\subseteq\mathcal J\subseteq\cI,
 \qquad
 \cI^{ar}\subseteq\mathcal J^r\subseteq\cI^r
 \quad(r\geq1).
\end{equation}
Such a shrinking may be made on $X$: for a proper map, every open
neighborhood of its compact fiber contains the inverse image of a
sufficiently small neighborhood of the image point.

The proper analytic direct-image theorem makes
$R^1\pi_*\mathcal E$ coherent.  The analytic theorem on formal
functions identifies its completed stalk with
\begin{equation}\label{analytic:formal-functions}
 \widehat{(R^1\pi_*\mathcal E)_x}
 \simeq
 \varprojlim_{r\geq1}H^1(Y,\mathcal E/\mathcal J^r\mathcal E).
\end{equation}
We use \cite[\S6, Haupts\"atze~I and~IIa,
p.~59]{Grauert1960}: the hypotheses are a proper holomorphic map
and a coherent analytic sheaf.  In particular no flatness of
$\pi$ is required.  The quotients in
\eqref{analytic:formal-functions} are supported on the infinitesimal
fibers over $x$, so their cohomology is independent of a further
shrinking of the representative.

For clarity, the right side vanishes by an explicit factorization,
without identifying the individual cohomology groups for the two
filtrations.  The inclusions in \eqref{analytic:cofinal} give
\[
 \mathcal J^{ar}\mathcal E
 \subseteq\cI^{ar}\mathcal E
 \subseteq\mathcal J^r\mathcal E.
\]
Hence the transition map from index $ar$ to index $r$ in
\eqref{analytic:formal-functions} factors through
\[
 H^1(Y,\mathcal E/\cI^{ar}\mathcal E)=0
\]
by Lemma~\ref{analytic:graded-vanishing}.  The inverse system is
therefore pro-zero, and its inverse limit is zero.  Since the stalk
$(R^1\pi_*\mathcal E)_x$ is a finite module over the Noetherian
local ring $\OO_{X,x}$, its zero completion implies that the stalk
itself is zero, by Nakayama's lemma.  Away from $x$ the map $\pi$
is an isomorphism, so the same higher direct image vanishes there.
We obtain $R^1\pi_*\mathcal E=0$.

The sheaf $\pi_*\mathcal E$ is also coherent.  Cartan's
Theorem~B on the Stein space $X$, together with the Leray spectral
sequence for $\pi$, now gives
\[
 H^1(Y,\mathcal E)=0.
\]
The cohomology sequence of
\[
 0\longrightarrow\mathcal E\longrightarrow\cF
 \longrightarrow\cF/\cI^9\cF\longrightarrow0
\]
therefore lifts both $\bar v_i$ to actual holomorphic sections $v_i$
of $\cF$.

It remains to check that arbitrary such lifts retain the required
frame.  In each divisor chart, let $v_i^{(8)}$ be the polynomial
representative of the prescribed jet.  It belongs to $\cI\cF$, and
\[
 v_i-v_i^{(8)}\in\cI^9\cF.
\]
Thus the change in their wedge belongs to
$\cI^{10}\det\cF$.  By Lemma~\ref{analytic:realization}, the
logarithmic determinant of the lifted pair is still
$t^4(1+O(t))$; its leading coefficient is $1$ even at the common
zeros of the leading row $r$.  Its ordinary determinant is
\begin{equation}\label{analytic:ordinary-determinant}
 v_1\wedge v_2=t^5 u(t,z)\,\partial_t\wedge\partial_z,
 \qquad u(0,z)=1.
\end{equation}
The unit condition holds on a neighborhood of each point of $C$.
Compactness gives one neighborhood on which all these local factors
are nonvanishing.  Properness allows a further shrinking of $X$
whose inverse image lies in this neighborhood.  Equation
\eqref{analytic:ordinary-determinant} then proves that $v_1,v_2$
frame $\cT_{Y\setminus C}$.
\end{proof}

\subsection{The singular point and its derivations}

\begin{proposition}\label{prop:analytic-counterexample}
There is a normal complex surface germ $(X,x)$ with an isolated
singular point and free analytic derivation sheaf of rank two:
\[
 \cT_X=\mathcal{H}om_{\OO_X}(\Omega_X^1,\OO_X)
 \simeq\OO_X^{\oplus2}.
\]
It has a resolution $\pi:Y\to X$ whose exceptional set is the
smooth curve $C$ of genus $21$ and self-intersection $-10$.
\end{proposition}

\begin{proof}
Take the contraction constructed in
Lemma~\ref{analytic:actual-lifting}.  Suppose first that $x$ were a
smooth point.  A sufficiently small neighborhood of $x$ would admit
a nowhere-vanishing holomorphic two-form $\omega$.  Its pullback
would be holomorphic on $Y$ and nonvanishing off $C$.  Thus
\[
 \operatorname{div}(\pi^*\omega)=aC\qquad(a\geq0),
 \quad\hbox{and hence}\quad K_Y\cdot C=aC^2\leq0.
\]
Adjunction, however, gives
\[
 K_Y\cdot C=2g(C)-2-C^2=40+10=50,
\]
a contradiction.  The point $x$ is singular, and it is the only
singular point after the chosen shrinking.

Write $X^\circ=X\setminus\{x\}$ and
$j:X^\circ\hookrightarrow X$.  The two fields of
Lemma~\ref{analytic:actual-lifting} descend to a holomorphic frame
$w_1,w_2$ on $X^\circ$.  Normality in dimension two gives the
Hartogs extension property~\cite[Chapter~II, Corollary~7.8]{Demailly2012}
\begin{equation}\label{analytic:hartogs}
 \OO_X\simeq j_*\OO_{X^\circ}.
\end{equation}
This property extends the fields as derivations, as follows.  Embed
a representative of $(X,x)$ in a smooth analytic germ with
coordinate functions $x_1,\ldots,x_e$.  For each $i$ and $\alpha$,
the holomorphic function $w_i(x_\alpha)$ on $X^\circ$ extends
uniquely to a function $a_{i\alpha}$ on $X$.  If $f$ is any local
defining equation of the embedding, then
\[
 \sum_{\alpha=1}^e
 a_{i\alpha}\frac{\partial f}{\partial x_\alpha}=0
\]
on $X^\circ$ and hence on $X$, by density and reducedness.  The
tuple $(a_{i\alpha})_\alpha$ consequently defines a section of
$\mathcal{H}om_{\OO_X}(\Omega_X^1,\OO_X)$ extending $w_i$.  We
continue to denote that extension by $w_i$.

For any local analytic derivation $w$ on $X$, the frame on
$X^\circ$ gives unique holomorphic coefficients $h_1,h_2$ there
such that $w=h_1w_1+h_2w_2$.  Both coefficients extend by
\eqref{analytic:hartogs}.  The extended equality holds on $X$,
since applying its difference to the embedding coordinates gives
holomorphic functions vanishing on the dense open set $X^\circ$.
This proves surjectivity of the map
\[
 \OO_X^{\oplus2}\longrightarrow\cT_X,
 \qquad(h_1,h_2)\longmapsto h_1w_1+h_2w_2.
\]
If a pair lies in its kernel, both coefficients vanish on
$X^\circ$ because $w_1,w_2$ are a frame there, and therefore vanish
on $X$.  The map is also injective, proving the assertion.
\end{proof}

\begin{remark}\label{analytic:invariants}
The construction also determines its canonical discrepancy and a
lower bound for its geometric genus.  The two-form $\omega^\circ$
dual to the frame $w_1,w_2$ generates the canonical sheaf on
$X^\circ$.  Extending reflexively across the normal surface point
gives $\omega_X\simeq\OO_X$; equivalently, we may choose the
canonical divisor $K_X=0$.  Since a normal surface is Cohen--Macaulay,
this reflexive canonical sheaf is its dualizing sheaf, so $X$ is
also Gorenstein; this surface criterion is recalled in
\cite[\S2.4, Definition~2.18 and the following paragraph]{FlennerZaidenberg2001}.
By
\eqref{analytic:ordinary-determinant}, its pullback is
\[
 \pi^*\omega^\circ=t^{-5}u(t,z)^{-1}\,dt\wedge dz.
\]
It is meromorphic on $Y$ with divisor exactly $-5C$.  Thus
\[
 K_Y=\pi^*K_X-5C,
\]
so the discrepancy of $C$ is $-5$.

For the geometric genus, put
$\OO_{nC}=\OO_Y/\cI^n$.  Riemann--Roch and the established
equality $h^0(C,L)=2$ give
\[
 h^1(C,L)=h^0(C,L^3)=12.
\]
In the sequence
\[
 0\longrightarrow L\longrightarrow\OO_{2C}
 \longrightarrow\OO_C\longrightarrow0,
\]
the constant section of $\OO_C$ lifts to $\OO_{2C}$.  Since
$H^2(C,L)=0$, its cohomology sequence gives
\[
 h^1(\OO_{2C})=h^1(C,L)+h^1(C,\OO_C)=12+21=33.
\]
For every $n\geq1$, the sequence with kernel $L^n$ similarly shows
that $H^1(\OO_{(n+1)C})\to H^1(\OO_{nC})$ is surjective.
The theorem on formal functions, with the cofinal filtrations in
\eqref{analytic:cofinal}, identifies their inverse limit with the
completion of $(R^1\pi_*\OO_Y)_x$.  This stalk has finite length,
since the coherent higher direct image is supported at $x$, and
thus equals its completion.  The inverse limit maps onto
$H^1(\OO_{2C})$, proving
\[
 p_g(X,x)=\dim_\CC (R^1\pi_*\OO_Y)_x\geq33.
\]
The exceptional curve has genus $g=21$ and its dual graph has first
Betti number $b=0$.  Consequently $p_g-g-b\geq12$.
\end{remark}
